% IEEE IoT-J: one paragraph, ~150--250 words, no citations, no displayed math.
\begin{abstract}
Q-learning routers for underwater sensor networks rank forwarders using held neighbour state, yet the cost of keeping that state current depends on how the state travels. We study this dependence with DA-HELLO, a configurable HELLO scheduler for a host router. A node sends a HELLO when it first decodes a never-recorded sender, when its candidate membership or void status changes, and, optionally, on a slow refresh; a local score-change test is also included. We evaluate it on two Q-learning routers, MARV and QELAR, each with state carried in data headers and in a constructed HELLO-only variant, in static, drifting, and node-loss scenarios. Comparisons are paired over 40 seeds, statements and margins were fixed before the runs, and idle listening is counted in energy. When state travels only in HELLOs, DA-HELLO delivers 2.1 to 3.0 percentage points more than a jittered 120~s periodic schedule on MARV, with 43 to 53\% fewer control bytes and 8 to 10\% lower energy per bit. On QELAR it uses 45 to 52\% fewer bytes and gains 5.0 points in the static scenario, but is 2.2 and 1.7 points lower in the drifting ones, so non-inferiority at 3 points is not shown. When state is carried in data headers, DA-HELLO matches sending no HELLOs, whereas periodic and Trickle schedules lose delivery. The score-change test shows no demonstrated benefit. The configuration differs by host class, and the results come from simulation of two re-implemented routers.
\end{abstract}
